\scanpage{63}
\mspage{50}
% archive sheet number: 62

\uline{Théorème} 3.1. Soit $(M,W,C)$ une catégorie à cofibrations (cf page 39), \inserted{$M$ stable par $\unsure{\varinjlim|\lim}$ finies.} Pour tout ensemble ordonné $I$, considérons le triple $(M(I^{\unsure{\circ|0}}),W_{I^\circ},C_{I^\circ})$ correspondant, où $C_{I^\circ}$ est défini au \S{}2 (cf à ce sujet le th.~2.1). Supposons que \deleted{\unsure{soit}} $I$ soit un ens ordonné \uline{artinien}, (p.\ ex. $I$ \unsure{loc.}$^{\unsure{t|mt}}$ fini i.e. les $I_{/i}$ finis). Supposons de plus qu'on soit dans l'un des deux cas suivants

a) L'espace $\unsure{\mathrm{Top}}(I)$ est noethérien i.e. toute suite croissante d'ouverts de $I$ est stationnaire (p.\ ex. $I$ fini).

b) $C$ et $TC=C\cap W$ sont stables, non seulement par composition \unsure{ordinaires|ordinaire} mais aussi par compositions transfinies. Dans ce cas, on suppose aussi $M$ stable par \inserted{petites} $\varinjlim$ \deleted{``quelconques''} filtrantes (donc par \inserted{petites} $\varinjlim$ quelconques).

Sous ces conditions, $(M(I^\circ),W_{I^\circ},C_{I^\circ})$ est également une catégorie à cofibrations.

\uline{\deleted{\unsure{Corollaire}}} \inserted{Prop\unsure{.}} 3.2. Soit $u:I\longrightarrow I'$ une application croissante entre ensembles ordonnés, d'où un couple de foncteurs adjoints

\diagram{\begin{tikzcd}[ampersand replacement=\&,column sep=large]
 (3.\unsure{2|1}.1) \& M(I^\circ)\arrow[r,shift left=0.6ex,"{u_!}"] \& M(I'^\circ)\arrow[l,shift left=0.6ex,"{u^*}"]
\end{tikzcd}}
